\documentclass[aspectratio=169,11pt]{beamer}

\usetheme{default}
\usecolortheme{dove}
\usepackage{amsmath}
\usepackage{booktabs}

\definecolor{accent}{RGB}{70,55,145}
\setbeamercolor{structure}{fg=accent}
\setbeamercolor{frametitle}{fg=accent}
\setbeamertemplate{navigation symbols}{}
\setbeamertemplate{footline}[frame number]

\newcommand{\FLS}{\mathrm{FLS}}

\title{Recovering FLS Sampling-Composition Shocks}
\subtitle{Entropy calibration using OEWS areas}
\author{H-2A Project}
\date{}

\begin{document}

\begin{frame}
  \titlepage
\end{frame}

\begin{frame}{Idea}
  \begin{itemize}
    \item The FLS reports one wage for each AEWR region and year.
    \item Realized respondents may be distributed unevenly within the region.
    \item This composition can move the FLS wage.
    \item Recover the smallest geographic reweighting that reproduces it.
  \end{itemize}

  \begin{block}{Interpretation}
    An implied sampling-composition shock, not the confidential FLS weights.
  \end{block}
\end{frame}

\begin{frame}{Objects}
  For region $r$, year $t$, and OEWS area $a$:

  \begin{itemize}
    \item $w_{art}$: OEWS wage for the AEWR occupations.
    \item $p_{art}$: prior weight from county farm employment.
    \item $\bar w^{\FLS}_{rt}$: regional FLS wage.
    \item $q_{art}$: calibrated weight.
  \end{itemize}

  OEWS areas are the wage-information units. Counties remain the units used to
  construct commuting zones and dissimilarity clusters.
\end{frame}

\begin{frame}{Wage calibration}
  \[
    \min_q \sum_a q_{art}\log\left(\frac{q_{art}}{p_{art}}\right)
  \]
  subject to
  \[
    \sum_a q_{art}=1,
    \qquad
    \sum_a q_{art}w_{art}=\bar w^{\FLS}_{rt},
    \qquad
    q_{art}\geq0.
  \]

  The solution exponentially tilts the prior:
  \[
    q_{art}\propto p_{art}\exp(\lambda_{rt}w_{art}).
  \]

  One wage moment does not identify literal sampling weights. KL divergence
  selects the closest implied distribution.
\end{frame}

\begin{frame}{Regularization path}
  Let
  \[
    \bar w^0_{rt}=\sum_a p_{art}w_{art}.
  \]

  Balance to
  \[
    \bar w^\gamma_{rt}
      =(1-\gamma)\bar w^0_{rt}+\gamma\bar w^{\FLS}_{rt},
    \qquad
    \gamma\in\{0,.25,.5,.75,1\}.
  \]

  \begin{itemize}
    \item $\gamma=0$: employment prior.
    \item $\gamma=1$: exact FLS match when feasible.
    \item Intermediate values show sensitivity to the composition assumption.
  \end{itemize}
\end{frame}

\begin{frame}{Dissimilarity IV}
  For target unit $k$, let $D(k)$ be its dissimilar donor cluster:
  \[
    Z^\gamma_{kt}
      =\frac{\sum_{a\in D(k)}q^\gamma_{ar,t-1}w_{ar,t-1}}
      {\sum_{a\in D(k)}q^\gamma_{ar,t-1}}.
  \]

  \begin{itemize}
    \item Drop an entire OEWS area if it touches the target cluster.
    \item Renormalize the remaining donor mass.
    \item Estimate levels with unit and year FE, or changes with year FE.
  \end{itemize}
\end{frame}

\begin{frame}{Why add moments?}
  One wage restriction can be satisfied by many geographic weight vectors.

  Additional moments can:
  \begin{itemize}
    \item reduce dependence on the prior;
    \item connect weights to known FLS sampling dimensions;
    \item distinguish composition from a common wage shock;
    \item provide overidentifying checks.
  \end{itemize}

  \begin{alertblock}{Requirement}
    Each moment needs both a regional FLS target and an area-level analogue.
  \end{alertblock}
\end{frame}

\begin{frame}{Feasible public moments}
  \small
  \begin{tabular}{lll}
    \toprule
    FLS target & Area analogue & Treatment \\
    \midrule
    150+ day worker share
      & Census of Agriculture Table 7
      & Soft balance \\
    Quarterly worker shares
      & QCEW monthly employment
      & Soft balance \\
    Field and livestock wages
      & OEWS occupation wages
      & Sensitivity \\
    Hours per worker
      & ACS PUMS usual hours
      & Validation \\
    \bottomrule
  \end{tabular}

  \vspace{0.8em}
  The first two moments most directly describe worker composition over the
  year.
\end{frame}

\begin{frame}{Worker-duration moment}
  FLS target:
  \[
    d^{\FLS}_{rt}
      =\frac{\sum_s N^{150+}_{rts}}{\sum_s N^{\mathrm{all}}_{rts}},
    \qquad s\in\{Jan,Apr,Jul,Oct\}.
  \]

  Area analogue:
  \begin{enumerate}
    \item Download county worker counts from Census of Agriculture Table 7.
    \item Use 2007, 2012, 2017, and 2022.
    \item Aggregate counts to OEWS areas.
    \item Interpolate area shares between census years.
  \end{enumerate}

  Use a soft constraint: the Census counts annual workers, while the FLS uses
  reference weeks.
\end{frame}

\begin{frame}{Seasonal-employment moment}
  FLS target:
  \[
    s^{\FLS}_{rtm}
      =\frac{N^{\FLS}_{rtm}}
      {\sum_hN^{\FLS}_{rth}}.
  \]

  Area analogue:
  \begin{enumerate}
    \item Download quarterly QCEW files.
    \item Keep January, April, July, and October employment.
    \item Use private NAICS 111 and 112; exclude agricultural services.
    \item Aggregate counties to OEWS areas and form within-year shares.
  \end{enumerate}

  Balance three quarterly shares. Use a soft penalty because QCEW omits many
  workers on small farms.
\end{frame}

\begin{frame}{Other candidates}
  \textbf{Field and livestock wages}
  \begin{itemize}
    \item FLS regional targets already exist in \texttt{fls\_region.parquet}.
    \item OEWS cores are SOC 45-2092 and 45-2093.
    \item Other AEWR occupations span both farm types, so the crosswalk is
      imperfect. Use only as a sensitivity check.
  \end{itemize}

  \textbf{Hours per worker}
  \begin{itemize}
    \item Use ACS PUMS \texttt{UHRSWORK} for wage workers in NAICS 111/112.
    \item PUMA geography and residence-based measurement make this a validation
      moment, not an exact constraint.
  \end{itemize}
\end{frame}

\begin{frame}{Prior variables are not balancing moments}
  Useful predictors without public regional FLS targets include:

  \begin{itemize}
    \item peak-worker size bins;
    \item gross-sales class;
    \item cropland acres and crop mix;
    \item livestock inventories and farm type.
  \end{itemize}

  Use them to improve $p_{art}$ or for national validation. Do not impose a
  balance restriction without a corresponding realized regional FLS target.
\end{frame}

\begin{frame}{Soft multi-moment calibration}
  For area variables $x_{ja}$ and FLS targets $\bar x_j^{\FLS}$:
  \[
    \min_q
      \sum_a q_a\log(q_a/p_a)
      +\sum_j\frac{\rho_j}{2}
      \left(
        \frac{\sum_aq_ax_{ja}-\bar x_j^{\FLS}}{s_j}
      \right)^2.
  \]

  \begin{itemize}
    \item $s_j$ standardizes the moments.
    \item $\rho_j$ controls how strongly each target is imposed.
    \item Report imbalance, KL divergence, effective area count, and maximum
      weight.
  \end{itemize}
\end{frame}

\begin{frame}{Recommended sequence}
  \begin{enumerate}
    \item Combined wage only: current benchmark.
    \item Add the 150-day worker share.
    \item Add quarterly worker shares.
    \item Test separate wage proxies.
    \item Use hours as validation.
  \end{enumerate}

  Compare weights, support, and first-stage strength after each block.

  \vspace{0.8em}
  \small
  Sources: \href{https://www.nass.usda.gov/Surveys/Guide_to_NASS_Surveys/Farm_Labor/}{FLS},
  \href{https://www.nass.usda.gov/Publications/AgCensus/2022/}{Census of Agriculture},
  \href{https://www.bls.gov/cew/}{QCEW}, and
  \href{https://www.census.gov/programs-surveys/acs/microdata.html}{ACS PUMS}.
\end{frame}

\end{document}
